\section{Screening Framework and Legal Scope}
\label{app:framework_v18}

The model used in the main text is intentionally spare. Its job
is not to micro-found cartel formation or side payments. Those
issues are treated in the cartel literature
\citep{asker2010leniency,marshall2012economics}. The narrower
job here is to show why a firm that repeatedly loses, while
continuing to appear in procurement records, can be a useful
screening object before bid-level proof is available.

\subsection{Primitives}
\label{app:framework_primitives}

Consider a procurement environment with repeated tender-items.
In a cartel-targeted tender, one firm is assigned the winning
role and submits the designated bid $b^*$. A cover bidder submits
a bid $b_\ell>b^*$ and is not meant to win. The cover role has a
cost $c_1>0$. Winning by mistake or by deviation creates an
internal cartel sanction $\kappa>0$, interpreted broadly as the
loss of cartel discipline, future rents, or capacity-allocation
rights. Let $r$ denote the private gain from deviating into the
winner role in a tender where the firm was supposed to provide
cover.

Firms outside the winning role are of two reduced-form types.
Type $C$ firms are cartel-deployed cover bidders. Type $G$ firms
are genuine losing participants: they may be inefficient,
capacity constrained, inexperienced, or simply unlucky. The
empirical screen conditions on the observed event
$W_i=0$, where $W_i$ is the firm's number of wins, and ranks
firms by participation count $T_i$.

\begin{lemma}[Zero-win cover role]
\label{lem:v18_zero_win}
If the internal sanction from winning in a cover role exceeds
the private deviation gain, $\kappa>r$, a cover bidder's
equilibrium bid lies above the designated winning bid:
$b_\ell>b^*$. Hence $\Pr(W_i>0\mid C)=0$ over tenders in which
the firm is deployed as cover.
\end{lemma}

\begin{proof}
If the cover bidder submits $b_\ell>b^*$, it loses and receives
the cover-role payoff net of the deployment cost. If it deviates
to a winning bid, it gains $r$ but pays the sanction $\kappa$.
The deviation is unprofitable when $\kappa>r$. The result fixes
the win probability of the cover role, not the exact numerical
cover bid. Any bid above the designated winner's bid is
consistent with the reduced-form cover role. 
\end{proof}

Lemma~\ref{lem:v18_zero_win} explains why the zero-win
conditioning set is natural. It does not say that all zero-win
firms are cartel participants. Many genuine firms also never
win. The screen's identifying content comes from the second
observable, participation intensity.

\subsection{Participation as a Ranking Statistic}
\label{app:framework_participation}

Within the zero-win set, let participation counts follow a
Poisson approximation:
\[
  T_i \mid \tau_i, q \sim \text{Poisson}(\lambda_q \tau_i),
  \qquad q\in\{C,G\},
\]
where $\tau_i$ is the length of time firm $i$ is observed and
$\lambda_q$ is the type-specific participation rate. The
screening assumption is
\[
  \lambda_C>\lambda_G
  \quad\text{on the conditioning set } W_i=0.
\]
This is the substantive assumption behind the frequent-loser
screen. Cover bidders are repeatedly deployed to provide losing
participation; ordinary zero-win firms participate less often.

\begin{proposition}[Monotone loser-side suspicion]
\label{prop:v18_monotone}
Under the Poisson approximation and $\lambda_C>\lambda_G$, the
posterior probability $\Pr(C\mid T_i=t,W_i=0,\tau_i)$ is
non-decreasing in $t$. Consequently, $T_i$ and
$\log(1+T_i)$ are rank-equivalent screening statistics for
cover-bidder exposure within the zero-win stratum.
\end{proposition}

\begin{proof}
By Bayes' rule, the posterior odds of type $C$ against type $G$
are proportional to the likelihood ratio
\[
  \frac{\Pr(T_i=t\mid C,\tau_i)}
       {\Pr(T_i=t\mid G,\tau_i)}
  =
  \exp[-(\lambda_C-\lambda_G)\tau_i]
  \left(\frac{\lambda_C}{\lambda_G}\right)^t .
\]
Because $\lambda_C>\lambda_G$, the likelihood ratio is increasing
in $t$. This is the monotone-likelihood-ratio property of the
Poisson family \citep{karlin1956mlr}. The posterior increases
with participation count, and the log transformation
$\log(1+T_i)$ preserves the ranking.
\end{proof}

The empirical binary rule used in the main text is a coarsening
of Proposition~\ref{prop:v18_monotone}. It flags firms above the
median plus $1.5$ times the interquartile range of zero-win
participation. The threshold is not a structural parameter. It
is an auditable way to turn a monotone ranking into an
operational list. Agencies could instead use the continuous
score, a capacity-constrained top-$k$ list, or a threshold
calibrated to their investigative budget.

\subsection{Why the Screen Is Not Proof}
\label{app:framework_scope}

The proposition delivers a ranking result, not an adjudication
result. Even if the model is correct, a high value of
$\log(1+T_i)$ raises the posterior probability of cover-bidder
exposure only relative to other zero-win firms. It does not
identify an agreement, assign cartel membership, prove that any
specific bid was cover, or establish a damages base. Those
claims require evidence from the richer record: bid values,
communications, leniency material, repeated winner rotation, or
other case-specific facts.

This boundary is the legal reason for the architecture in the
main text. The award layer is allowed to create suspicion because
it ranks firms on a statistic that is cheap, observable, and
monotone under a transparent behavioral assumption. The bid layer
remains the proof layer because liability turns on conduct and
agreement, not on suspicious participation alone.

\begin{proposition}[Legal-scope separation]
\label{prop:v18_scope}
Under Lemma~\ref{lem:v18_zero_win} and
Proposition~\ref{prop:v18_monotone}, the frequent-loser statistic
is sufficient for triage only. It is not sufficient for liability
or damages.
\end{proposition}

\begin{proof}
Triage requires an ordering of expected investigative value.
Proposition~\ref{prop:v18_monotone} supplies that ordering
within the zero-win stratum. Liability and damages require
different predicates: proof of agreement, conduct in particular
tenders, and an economic measure of harm. The statistic contains
participation and win information only. It omits the facts needed
for those predicates. The same statistic can be useful for
allocating forensic attention and insufficient for legal
adjudication.
\end{proof}
